\documentclass[11pt]{article}
\usepackage[a4paper,margin=22mm]{geometry}
\usepackage{fontspec}
\setmainfont{LinLibertine_R.otf}[BoldFont=LinLibertine_RB.otf,ItalicFont=LinLibertine_RI.otf,BoldItalicFont=LinLibertine_RBI.otf]
\usepackage{xcolor}
\usepackage{hyperref}
\hypersetup{colorlinks=true,urlcolor=blue}
\setlength{\parindent}{0pt}
\setlength{\parskip}{6pt}
\sloppy
\emergencystretch=2em
\begin{document}
\section*{Historical Britannica: extracted source}
Encyclopaedia Britannica - Seventh edition, Volume 7, CLO-Dialling - EB.15\par
{\small Source row 1843. Parser: nvidia/nemotron-3.5-parse.\par}
{\small This is model-extracted text, not a corrected transcription. OCR errors, reading order and region alignment require human review.\par}
\hrule\medskip
{\small\itshape CONIC SECTIONS.\par}

{\scriptsize\color{gray} Region region-0; provider order 0.\par}
\medskip

223
\par

{\scriptsize\color{gray} Region region-1; provider order 1.\par}
\medskip

angle HKp is equal to Hp: But HK being equal to Hà, for each has been proved equal to HP, the angles HKA. Hkk are equal (5, 1, E), and adding to these the right angle PKK, piks; the angles HKP. Haynes equals but those have been proved equal to HP. Hp: therefore there last are equal, and the line HF makes equal angles with FP-lp.
\par

{\scriptsize\color{gray} Region region-2; provider order 2.\par}
\medskip

Con. 1. Perpendication drawn from the intersection of two tangente, to the draw down from the points of contact through the focus are equal. For HL: Hi, being drawn perpendicular to HP, PF, the triangles HFI, HFG, are manifolding equal (9). L.), where the Hl = HL.
\par

{\scriptsize\color{gray} Region region-3; provider order 3.\par}
\medskip

Con. 2. Perpendication from the intersection of two tangents to diameters passing through the points of con› tace equal.
\par

{\scriptsize\color{gray} Region region-4; provider order 4.\par}
\medskip

Draw Ghq through HP perpendicular to PK, ph, and be› ceus the triangular HPG; HP) have HP common to both the angles at P, equal, and the angles of and I right angles, the triangles are in every way equal (28, 1, E.), and hence HG is equal to HL. In like manner it is proved that Hp a equal to Hr; but HI is equal to Hd, therefore HG is equal to Hg.
\par

{\scriptsize\color{gray} Region region-5; provider order 5.\par}
\medskip

{\small\itshape Pero. VI.\par}

{\scriptsize\color{gray} Region region-6; provider order 6.\par}
\medskip

If a straight line be drawn from the intersection of two tangents to the focus, and another perpendicular to the direction; these will make equal angles with the tangents.
\par

{\scriptsize\color{gray} Region region-7; provider order 7.\par}
\medskip

C
H
C
C
K
E
I
I
D
P
P
N
\par

{\scriptsize\color{gray} Region region-8; provider order 8.\par}
\medskip

Let F be the focus of a parabola, and K4 the directrix; and let straight lines HP, Hp, which intersect each other at H, touch the parabola or P and , also let HP be drawn to the focus, and HE perpendicular to the directey, the angles PHP, pHP are equal.
\par

{\scriptsize\color{gray} Region region-9; provider order 9.\par}
\medskip

The same construction being made as in Prop. V.
\par

{\scriptsize\color{gray} Region region-10; provider order 10.\par}
\medskip

In the triangles HEK. HEÀ it may be shown, as in that proportion, the HK is equal to Hk, and therefore that the angle HKE is equal to the angle HRT (A, 1, E.). The angles HKE,HE are also equal; therefore the angles KEH e one equal (26, 1, E.). Now the angle KHE=KHP + PHE; but the triangles KHP, PHP are in every way and (as was shown in Prop. V.) Therefore KHP=(HP, and hence
\par

{\scriptsize\color{gray} Region region-11; provider order 11.\par}
\medskip

KE=FHP + PHE= PHE+ SPHE. In the same way it appears that
\par

{\scriptsize\color{gray} Region region-12; provider order 12.\par}
\medskip

KE= FHP+ pHt= FHE+SPHp; therefore FHE + SPHE= FHE+ EPHp; and hence ZPHE=3pHT and PHE= pHF.
\par

{\scriptsize\color{gray} Region region-13; provider order 13.\par}
\medskip

{\small\itshape Pero. VII.\par}

{\scriptsize\color{gray} Region region-14; provider order 14.\par}
\medskip

If two tangents to a parabola be at the extremities of a chord, and a third tangent be parallel to the chord ; the
\par

{\scriptsize\color{gray} Region region-15; provider order 15.\par}
\medskip

part of this tangent intercepted between the other two Parabola, is bisected at the point of contact.
\par

{\scriptsize\color{gray} Region region-16; provider order 16.\par}
\medskip

Z
H
Z
Z
Z
Z
M
K
N
Z
Z
T
E
I
\par

{\scriptsize\color{gray} Region region-17; provider order 17.\par}
\medskip

Les HD, Hd, be tangents at the extremities of the chord Dd and KRB a tangent parallel to Dd, meeting the other tangents in K and Y, the line K4 is bisected at P, the point of contact.
\par

{\scriptsize\color{gray} Region region-18; provider order 18.\par}
\medskip

From Hi, K, k, the intersections of the tangents, dray perpendicules to the dimeters passing through their points of contact; viz. H. H, perpendicular to DL and dl, and KN, KN, perpendicular to DL and PE, and km, km, perpendicular to di and PE.
\par

{\scriptsize\color{gray} Region region-19; provider order 19.\par}
\medskip

The triangle SHI, DKN, are manifestly equiangular, also the triangle dHf, day : therefore
\par

{\scriptsize\color{gray} Region region-20; provider order 20.\par}
\medskip

HD: DK= IH: KM (4, 6, E.), and Hdt: dk= His .m.
\par

{\scriptsize\color{gray} Region region-21; provider order 21.\par}
\medskip

But because K4 is parallel to Dd,
\par

{\scriptsize\color{gray} Region region-22; provider order 22.\par}
\medskip

HD: DK= Ht/zk (2, 6, E.); therefore Ht KM= Hin.
\par

{\scriptsize\color{gray} Region region-23; provider order 23.\par}
\medskip

Now HI= Hi (2 Cor. 5), therefore KM= km.
\par

{\scriptsize\color{gray} Region region-24; provider order 24.\par}
\medskip

But KM= NV, and Km= Kon (2 Cor 5);
\par

{\scriptsize\color{gray} Region region-25; provider order 25.\par}
\medskip

therefore KN= kn.
\par

{\scriptsize\color{gray} Region region-26; provider order 26.\par}
\medskip

And since KN:kn= KP:(CP. 6, E.), therefore KP= kl:
\par

{\scriptsize\color{gray} Region region-27; provider order 27.\par}
\medskip

{\small\itshape LemMA\par}

{\scriptsize\color{gray} Region region-28; provider order 28.\par}
\medskip

Les KL be a triangle, having its base Ll bisected at p, and let Hå, any straight line parallel to the base, and terminates by the sides. As bisected at P, then Y, p., the points of insertion, and K, the vertex of the triangle, are in the same straight lines; and that line bisects Dd, any other straight line parallel to the base.
\par

{\scriptsize\color{gray} Region region-29; provider order 29.\par}
\medskip

H
Z
Z
H
I
D
B
\par

{\scriptsize\color{gray} Region region-30; provider order 30.\par}
\medskip

Complete the parallelograms KHPM, KlpN. The triangles KHÅ, KLU being similar, and Hå, L similarly divided at J and .
\par

{\scriptsize\color{gray} Region region-31; provider order 31.\par}
\medskip

KL: KL= Hä: LI= HP:lop, hence the parallelograms KHPM, KlpN are similar. Now they have a common angle at K; therefore they are about the same diameter, that is the points K, P, p are in the same straight line (20, 6, E.).
\par

{\scriptsize\color{gray} Region region-32; provider order 32.\par}
\medskip

Next, let Dd need Ky in R, then If DE (c KPt:ED)= PA: Ed; therefore DE is equal to Ed.
\par

{\scriptsize\color{gray} Region region-33; provider order 33.\par}
\medskip

{\small\itshape Prov. VIII.\par}

{\scriptsize\color{gray} Region region-34; provider order 34.\par}
\medskip

Any chord parallel to a tangent is bisected by the diame› ter which passes through the point of contact, or is an ordinate to that diameter.
\par

{\scriptsize\color{gray} Region region-35; provider order 35.\par}
\medskip
\newpage\section*{Source record}
bri-1843-f53bf32c975b\par
Image SHA-256: \texttt{f53bf32c975bfa37eca80f791ffd0a10}\par\texttt{1655e8736fe3ad5096fa5d028ca659ab}\par
Dataset revision: \texttt{309f765b1ecc83e0a7f7}\texttt{762e44b4a927a7555443}\par
National Library of Scotland, Encyclopaedia Britannica Illustrated (CC0).\par
\url{https://huggingface.co/datasets/NationalLibraryOfScotland/encyclopaedia_britannica_illustrated/viewer/default/train?row=1843}
\end{document}
